\section{Tổng quan}\label{sec:1}

\begin{keypoint}
\sig{gemm_accel_top} là bộ tăng tốc nhân ma trận INT8 gắn vào SoC Chipyard qua hai cổng AXI: một
AXI4-Lite slave cho CPU cấu hình, một AXI4-Full master để tự đọc/ghi DRAM. CPU dựng một
\emph{chuỗi descriptor} trong DRAM, ghi \sig{DESC_PTR}, ghi START, rồi chờ bằng polling. Accelerator
tự đọc từng descriptor, nạp A và B vào scratchpad, tính $C = A \times B$ trên mảng systolic 8×8 PE,
hậu xử lý tùy descriptor, rồi ghi kết quả ra DRAM. Thiết kế gồm đúng sáu khối chức năng, một miền
clock \sig{clk} 100 MHz, scratchpad ping-pong hai bank để nạp, tính và ghi của các descriptor kế
nhau chạy chồng nhau.
\end{keypoint}

\subsection{Phạm vi chức năng}\label{sec:1.1}
\begin{itemize}
\item Mỗi tác vụ là một chuỗi descriptor. Mỗi descriptor tính một phép GEMM: A là M×K, B là K×N, kiểu INT8 có dấu; tích lũy INT32; M ≤ 32, N ≤ 32, K ≤ 32.
\item Kết quả ra DRAM ở một trong hai dạng: C (INT32, không hậu xử lý) hoặc Y (INT8, sau bias, requantize, ReLU tùy descriptor, kẹp). Dạng nào do bit \sig{POST_EN} của descriptor quyết định.
\item CPU dựng descriptor trong DRAM, ghi \sig{DESC_PTR}, ghi START, rồi chờ kết quả bằng polling. DMA đọc/ghi DRAM qua AXI4-Full, hai kênh độc lập.
\item Nạp, tính và ghi của các descriptor kế nhau chạy chồng trên scratchpad ping-pong.
\item Dataflow của systolic array là \textbf{Output-Stationary}: PE(r, c) giữ bộ cộng dồn của $C[i_0+r][j_0+c]$ suốt k = 0 … K−1; A chảy sang đông, B chảy xuống nam, mỗi chu kỳ một PE; tile xong thì xả kết quả theo hàng sang VPU (D-30).
\item Ngoài phạm vi (không hỗ trợ): ngắt, $C \mathrel{+}= A \times B$, ma trận lớn hơn 32 trong một descriptor, các dataflow khác (Weight-Stationary, Input-Stationary), bộ nhớ ảo, GEMM lập trình trực tiếp bằng thanh ghi MMIO.
\end{itemize}

\begin{notebox}[Quyết định baseline D-01]
Hai nguồn khác nhau ở mô hình lập trình: \sig{spec_1} chỉ có chế độ chuỗi descriptor và dự trữ
offset 0x08…0x24, còn \sig{spec_2} còn giữ thêm một chế độ \emph{single} lập trình bằng các thanh ghi
\sig{MATRIX_DIM}, \sig{BASE_ADDR_IN}, \sig{BASE_ADDR_WEIGHT}, \sig{BASE_ADDR_OUT},
\sig{POST_CTRL}, \sig{BIAS_ADDR}, \sig{QUANT_M0}, \sig{QUANT_SHIFT} tại các offset đó, cùng
\sig{CTRL[0]} = START, \sig{CTRL[2]} = START\_CHAIN và \sig{STATUS[3]} = MODE.
Baseline chọn \sig{spec_1}: \textbf{chỉ có chế độ chuỗi descriptor}. Lý do: kiến trúc đã chốt là
descriptor-based DMA; chế độ register-programmed là một nhánh thiết kế thứ hai làm đúng cùng một
việc, nhân đôi đường cấu hình trong Register File và Main FSM (thêm tám bus cấu hình, mỗi bus ứng với một
thanh ghi RW, cùng tín hiệu \sig{start_kind} trên giao diện RF→FSM) và buộc STATUS mang thêm một bit chế độ. Tác động: 0x08…0x24 là Reserved; \sig{CTRL[2]} là
bit START duy nhất; \sig{STATUS[3]} dự trữ, đọc 0; không có tín hiệu \sig{start_kind}.
\end{notebox}

\subsection{Quy ước}\label{sec:1.2}
\begin{itemize}
\item M, N, K: A là M×K, B là K×N, C và Y là M×N. P: kích thước array vật lý (P×P PE), chốt bằng 8. \emph{word} = 32 bit; \emph{beat} = một lần truyền AXI 32 bit; \emph{tile} = khối P×P của C mà array tính trong một lượt.
\item \emph{Tác vụ}: từ khi START được chấp nhận đến khi STATUS báo DONE hoặc ERROR. \emph{Job}: một descriptor trong chuỗi. Tham số của job nằm trong descriptor. \sig{DESC_PTR} được chụp lúc START được chấp nhận.
\item \emph{Pha}: một bước của job — FETCH, kiểm tra, LOAD\_BIAS, LOAD\_A, LOAD\_B, COMPUTE, STORE\_C.
\item Số hex viết dạng 0x6000\_0000; số nhị phân dạng 4'b1111. Thời gian quy đổi theo chu kỳ \sig{clk} 10,0 ns.
\item Chiều tín hiệu ở biên module nội bộ: \sig{i_*} là input, \sig{o_*} là output. Tín hiệu AXI giữ nguyên tên \sig{s_axi_*} và \sig{m_axi_*}.
\end{itemize}

\subsection{Kiến trúc tổng thể}\label{sec:1.3}

Dữ liệu đi một chiều: DRAM → DMA → scratchpad → array → VPU → scratchpad → DMA → DRAM. Điều khiển đi
vòng: CPU → Register File → Main FSM → DMA và array, rồi về Main FSM bằng các xung done/error.
\sig{gemm_accel_top} không chứa logic dữ liệu và không chứa FSM; nó chỉ đấu nối sáu instance và đặt
bộ đồng bộ reset hai flip-flop.

Kiến trúc được vẽ thành hai hình, mỗi hình trên một trang nằm ngang riêng. Hình~\ref{fig:1-1} là
luồng dữ liệu: DMA, scratchpad, array và VPU cùng các giao diện dữ liệu giữa chúng; các cổng điều khiển
của bốn khối này chỉ vẽ thành bó dây ngắn ghi nhóm tên port, đi tới \sig{u_main_fsm}. Hình~\ref{fig:1-2}
là luồng điều khiển: Register File, Main FSM, các lệnh và trạng thái giữa FSM với bốn khối dữ liệu, net
\sig{soft_clr} và bộ đồng bộ reset; cổng dữ liệu của bốn khối không vẽ lại ở hình này. Mỗi dây là một
bó tín hiệu, ghi nhóm tên port ở hai đầu; danh sách đầy đủ thành viên và port map nằm trong model
(\sig{<desc id="model">}) của từng file SVG.

\diagfigwide{gemm_accel_top_dataflow}{Luồng dữ liệu của \sigc{gemm_accel_top}: DRAM → DMA → scratchpad
(A, B, BIAS) → array → VPU → scratchpad (C) → DMA → DRAM}{fig:1-1}

\diagfigwide{gemm_accel_top_control}{Luồng điều khiển của \sigc{gemm_accel_top}: CPU → Register File →
Main FSM → bốn khối dữ liệu, net \sigc{soft_clr} và bộ đồng bộ reset hai flip-flop}{fig:1-2}

\begin{spectable}{|L{0.2}|L{0.8}|}{Luồng dữ liệu và luồng điều khiển qua sáu khối}{tab:1-1}{\textbf{Chặng} & \textbf{Nội dung}}
CPU → RF & Ghi \sig{DESC_PTR} và \sig{CTRL[2]} qua AXI4-Lite; đọc \sig{STATUS}, \sig{PERF_*}, \sig{ERR_DESC_ADDR}, \sig{DESC_DONE_CNT} \\\hline

RF → FSM & \sig{o_start_pulse}, \sig{o_soft_req}, \sig{o_cfg_desc_ptr} \\\hline

FSM → DMA & Lệnh kênh đọc (\sig{o_rd_*}: FETCH, LOAD\_BIAS, LOAD\_A, LOAD\_B) và lệnh kênh ghi (\sig{o_wr_*}: STORE\_C) \\\hline

DRAM → DMA → SRAM & Kênh đọc AXI nhận beat 32 bit, ghi vào bank của buffer A, B hoặc BIAS. Descriptor không vào scratchpad: DMA trả 16 word về FSM qua \sig{o_desc_valid}, \sig{o_desc_idx}, \sig{o_desc_data} \\\hline

SRAM → Array & Array đọc vector A (\sig{o_a_rd_*}) và B (\sig{o_b_rd_*}), mỗi vector 64 bit, dữ liệu về sau đúng một chu kỳ \\\hline

Array → VPU & Luồng hàng kết quả: \sig{o_row_valid}, \sig{o_row_data} 256 bit, \sig{o_row_idx}, \sig{o_row_col}, \sig{o_row_mask}. Không có đường ready \\\hline

VPU → SRAM & VPU đọc bias (\sig{o_bias_rd_*}) và ghi bank C (\sig{o_c_wr_*}, 256 bit, có mask) \\\hline

SRAM → DMA → DRAM & Kênh ghi đọc bank C qua \sig{o_spm_rd_*}, đẩy qua WFIFO rồi phát beat W ra AXI \\\hline

FSM → RF & \sig{o_done_pulse}, \sig{o_err_pulse}, \sig{o_err_code}, \sig{o_perf_*}, \sig{o_soft_busy}, \sig{o_err_desc_addr}, \sig{o_desc_done_cnt} \\\hline

FSM → mọi khối & \sig{o_soft_clr}: xung xóa đồng bộ một chu kỳ tới RF, DMA, SRAM, Array, VPU \\\hline

\end{spectable}

\subsection{Module}\label{sec:1.4}

\begin{spectable}{|L{0.25}|L{0.15}|L{0.6}|}{Module trong \sigc{gemm_accel_top}}{tab:1-2}{\textbf{Module} & \textbf{Instance} & \textbf{Vai trò}}
\sig{gemm_register_file} & \sig{u_regfile} & AXI4-Lite slave; \sig{DESC_PTR}, trạng thái, bộ đếm; chấp nhận START và SOFT\_RESET \\\hline

\sig{gemm_main_fsm} & \sig{u_main_fsm} & Bộ tuần tự duy nhất: kiểm tra descriptor; lập lịch ba engine Đọc, Tính, Ghi; ra lệnh cho DMA và array; báo DONE/ERROR; điều phối soft reset \\\hline

\sig{gemm_dma_engine} & \sig{u_dma} & AXI4-Full master, hai kênh độc lập; FETCH descriptor, LOAD\_A, LOAD\_B, LOAD\_BIAS, STORE\_C \\\hline

\sig{gemm_local_sram} & \sig{u_sram} & Bốn buffer A, B, C, BIAS; mỗi buffer hai bank ping-pong; mỗi bank chia sub-bank một cổng cho array và VPU \\\hline

\sig{gemm_systolic_array} & \sig{u_array} & Mảng 8×8 PE (\sig{gemm_pe}) output-stationary, làm lệch dữ liệu, bộ tuần tự tile, xả kết quả theo hàng \\\hline

\sig{gemm_vpu} & \sig{u_vpu} & Hậu xử lý trên đường xả kết quả: bias, requantize, ReLU, kẹp INT8; đi thẳng khi \sig{POST_EN} = 0 \\\hline

\end{spectable}

Ngoài sáu module trên, \sig{gemm_accel_top} chỉ có một bộ đồng bộ reset 2 flip-flop. Không có khối
thứ bảy, không có logic dữ liệu ở mức top.

\subsection{Tham số hệ thống}\label{sec:1.5}

\begin{spectable}{|L{0.3}|L{0.22}|L{0.48}|}{Tham số hệ thống}{tab:1-3}{\textbf{Tham số} & \textbf{Giá trị} & \textbf{Ghi chú}}
\sig{PE_ARRAY_SIZE} (P) & 8 & Giá trị hợp lệ duy nhất; giá trị khác bị chặn lúc elaboration \\\hline

Số PE & 64 & P × P \\\hline

Q & 2 & Số sub-bank của buffer B; dẫn xuất P/4 \\\hline

LP & 3 & $\log_2 P$ \\\hline

\sig{AB_VEC_WIDTH} & 64 bit & Vector A và B đọc từ SRAM; dẫn xuất 8·P \\\hline

\sig{ROW_VEC_WIDTH} & 256 bit & Hàng kết quả array → VPU → bank C; dẫn xuất 32·P \\\hline

\sig{MAX_DIM} & 32 & Giới hạn M, N, K \\\hline

\sig{AXI_DATA_WIDTH} & 32 & Cả AXI4-Lite và AXI4-Full \\\hline

\sig{MAX_BURST_BEATS} & 16 & Baseline chốt 16 (D-36); miền tham số 1 … 256 \\\hline

\sig{WFIFO_DEPTH} & 4 & FIFO ghi của DMA, 4 × 32 bit (localparam) \\\hline

\sig{MUL_STAGES} & 1 & Số tầng thanh ghi của bộ nhân requantize trong VPU \\\hline

\sig{VPU_LATENCY} (LAT) & 4 & $3 + \text{\sig{MUL_STAGES}}$; array dùng cùng giá trị làm \sig{L_VPU} \\\hline

NB & 2 & Số bank mỗi buffer (ping-pong) \\\hline

Buffer A / B / C / BIAS & 256 / 256 / 1024 / 32 word mỗi bank & Hai bank mỗi buffer; 1568 word mỗi bank; tổng 3136 word = 12,25 KiB \\\hline

\sig{SRAM_RD_LAT} & 1 chu kỳ & Cổng phía array, VPU và DMA \\\hline

\sig{MMIO_BASE} & 0x6000\_0000, vùng 4 KiB & \\\hline

\sig{DRAM_BASE} & 0x8000\_0000 & \\\hline

Burst outstanding & 1 đọc + 1 ghi & Mỗi kênh tối đa một burst chưa kết thúc \\\hline

\sig{TARGET_CLK_FREQ} & 100 MHz & Mục tiêu tổng hợp; chu kỳ 10,0 ns \\\hline

\end{spectable}

\begin{notebox}[Quyết định baseline D-04]
\sig{spec_1} cho \sig{PE_ARRAY_SIZE} nhận 4 hoặc 8 và khẳng định kết quả không phụ thuộc P;
\sig{spec_2} chốt P = 8 cố định. Baseline chốt \textbf{P = 8}, giá trị khác bị chặn lúc elaboration.
Lý do: kiến trúc đã chốt là PE 8×8, và P = 4 là một nhánh thiết kế thứ hai kéo theo hai bộ độ rộng
bus (32 bit và 128 bit cho vector A/B, 128 bit và 256 bit cho hàng kết quả), hai lịch skew và hai
công thức timing trong bốn module. Tác động: Q, LP, \sig{AB_VEC_WIDTH}, \sig{ROW_VEC_WIDTH} trở
thành hằng số; \sig{HW_CONFIG[7:0]} luôn đọc ra 8.
\end{notebox}

\begin{notebox}[Quyết định baseline D-05]
\sig{spec_1} cho \sig{MUL_STAGES} nhận 1 hoặc 2 và đặt \sig{L_VPU} của array mặc định 0 trong khi
vẫn yêu cầu \sig{L_VPU} = LAT — hai câu này mâu thuẫn nhau. \sig{spec_2} chốt \sig{MUL_STAGES} = 1,
LAT = 4. Baseline chốt \textbf{\sig{MUL_STAGES} = 1, LAT = \sig{L_VPU} = 4}, lấy
\sig{gemm_pkg::VPU_LATENCY} làm nguồn duy nhất; array và VPU đều kiểm giá trị này lúc elaboration.
Lý do: ở 100 MHz (10 ns) một tầng cho phép nhân 32×32 → 64 bit; \sig{MUL_STAGES} = 2 chỉ cần khi
đẩy tần số cao hơn. Chọn 1 giữ LAT nhỏ nhất, giảm chiều sâu đường ống VPU và khớp mặc định của
\sig{spec_2}. Nếu tổng hợp không đạt 100 MHz, cách sửa đã định trước là đặt \sig{MUL_STAGES} = 2
(LAT = \sig{L_VPU} = 5): đây là thay đổi tham số, không phải thay đổi kiến trúc.
\end{notebox}

\subsection{Clock và reset}\label{sec:1.6}

\begin{spectable}{|L{0.26}|L{0.74}|}{Clock và reset của hệ thống}{tab:1-4}{\textbf{Hạng mục} & \textbf{Quy định}}
Miền clock & Đúng một miền, \sig{clk}. Mọi flip-flop cập nhật tại cạnh lên. Không có CDC, không có clock gating, không có clock chia \\\hline

Tần số mục tiêu & 100 MHz, chu kỳ danh định 10,0 ns, duty 50 \% \\\hline

Đường tổ hợp dài nhất đã biết & Bộ nhân 32×32 → 64 bit của VPU (\sig{MUL_STAGES} = 1) và MAC INT8 một chu kỳ của \sig{gemm_pe} \\\hline

Reset ngoài & \sig{rst_n}, tích cực mức thấp, kích hoạt bất đồng bộ \\\hline

Đồng bộ reset & \sig{gemm_accel_top} có bộ đồng bộ 2 flip-flop: bật bất đồng bộ, nhả đồng bộ với \sig{clk}. Đầu ra \sig{rst_n_sync} cấp cho cả sáu module \\\hline

Độ trễ nhả reset & \sig{rst_n_sync} lên cao sau đúng 2 cạnh lên của \sig{clk} kể từ khi \sig{rst_n} lên cao \\\hline

Soft reset & \sig{soft_clr}: xung xóa \textbf{đồng bộ} một chu kỳ do Main FSM phát. Không đi vào chân reset của bất kỳ flip-flop nào \\\hline

Nội dung scratchpad & Không bị xóa bởi \sig{rst_n} cũng không bởi \sig{soft_clr}; coi là không xác định \\\hline

\end{spectable}

\begin{notebox}[Quyết định baseline D-20, D-21, D-22]
Không nguồn nào nêu chu kỳ clock; kiến trúc chốt đặt 100 MHz, nên baseline lấy chu kỳ danh định
10,0 ns và quy đổi mọi số đo thời gian theo đó (D-20). Nguồn nêu bộ đồng bộ 2 flip-flop nhưng không
nêu giá trị khởi tạo: baseline chọn \textbf{cả hai flip-flop có reset bất đồng bộ về 0}, nên
\sig{rst_n_sync} xuống thấp ngay khi \sig{rst_n} xuống, không cần clock chạy — bảo đảm accelerator
vào trạng thái reset khi SoC chưa có clock ổn định (D-21). Quy tắc hiện thực kèm theo: flip-flop
trạng thái dùng \sig{always_ff @(posedge clk or negedge rst_n_sync)}; flip-flop chỉ chứa dữ liệu
(\sig{acc} và \sig{res} của PE, nội dung RAM, thanh ghi dữ liệu trong đường ống VPU) dùng
\sig{always_ff @(posedge clk)} và không có giá trị reset. Nội dung scratchpad sau reset là
\textbf{không xác định}, không có mạch xóa (D-22); firmware không được đọc vùng đích trước DONE.
\end{notebox}
